\chapter{Folk: One Guitar, a Whole World}
\label{ch:genre-folk-deep}

\classquote{心之忧矣，我歌且谣。}{Book of Songs, ``Peach Trees in the Garden''}{My heart is grieved; I sing and I hum. Not for a stage, but to let the words out.}

\section{What Folk Is: Putting a Person Back in the Song}
\label{sec:folk-what}

Folk is easily mistaken for ``pop with a smaller budget'': one acoustic guitar, one singer, a few chords, anyone can do it. Good folk is not ``using fewer tracks''; it is \textbf{putting a person back in the song}---a specific person, in a specific place, telling a specific thing in a specific voice.

Its core listening experience usually has three layers:

\begin{longtable}{p{0.16\textwidth}p{0.36\textwidth}p{0.36\textwidth}}
\toprule
What the ear notices & What the listener hears first & What it means for arrangement \\
\midrule
\endfirsthead
\toprule
What the ear notices & What the listener hears first & What it means for arrangement \\
\midrule
\endhead
Voice in front & Someone talking in a room; breath, pauses, accent all present & The voice is the lead; the guitar is a second voice, not a wall \\
Clear texture & The grain of the acoustic guitar, fingers on strings, room noise & Few tracks, close capture, do not blur the midrange \\
Emotional rise and fall & Verses hold back; the chorus dares to open & Lift the chorus with another voice or instrument, not a fader \\
A story you can retell & You remember a place, an action, a person & Arrangement serves the story, not a ``vintage filter'' \\
\bottomrule
\end{longtable}

Folk can be rough or refined; it can be one guitar or strings, harmonica, accordion, light drums. The test is not track count but \textbf{whether each layer helps that person speak}. If removing a layer only makes the song quieter and the story unchanged, it can probably go; if a line loses its landing place, it should stay.

\section{The Route: From Folk Song Collecting to a Lonely Voice in the City}
\label{sec:folk-history}

\subsection{1940s--1960s: Folk Revival and ``One Person's Song''}
\label{sec:folk-history-1940s}

Modern folk has more than one route. One is the old folk-song tradition: work songs, field songs, narrative ballads, drinking songs---never meant for a concert hall. Another is the mid-20th-century English-language folk revival: Woody Guthrie wrote drifting, dust-bowl, unemployment, and land into songs in the 1940s; Pete Seeger collected, adapted, and performed folk songs, pulling ``folk'' out of the museum and back into the room.{\textsuperscript{[1]}}

In the 1950s, groups such as The Weavers brought folk onto commercial radio; in the 1960s, Bob Dylan pushed acoustic guitar, harmonica, and an ``unpolished'' voice to the front. The 1963 ``Blowin' in the Wind'' does not win through complex arrangement; one person, one guitar, one repeated question, and the listener supplies the answer.{\textsuperscript{[2]}}

\subsection{1970s: Singer-Songwriters Put Private Life in the Song}
\label{sec:folk-history-1970s}

In the 1970s, Joni Mitchell, James Taylor, Carole King, and other singer-songwriters put ``I'' into the song: love, city, travel, growing up, without needing a grand narrative. Joni Mitchell's 1971 album \emph{Blue} uses almost no fussy arrangement; voice, acoustic guitar, Appalachian dulcimer, and a little drums keep the vulnerability in front.{\textsuperscript{[3]}}

The period's inheritance is ``arrangement can subtract'': not every song needs a full band; one distinctive timbre---say, the metallic ring of a dulcimer---can be more identifiable than ten backing tracks.

\subsection{1990s: Chinese Campus Folk}
\label{sec:folk-history-china-1990s}

A key node in contemporary Chinese folk is 1994. Dadi Records released \emph{Campus Folk I}; Lao Lang sang ``My Deskmate'' and ``My Roommate upstairs,'' while Gao Xiaosong, Shen Qing, and others wrote youth, parting, dormitories, and playgrounds into songs.{\textsuperscript{[4]}} It turned folk from underground tapes and lawn sing-alongs into a mainstream sound heard on radio and television.

The success also left a side effect: many people assumed folk meant ``acoustic guitar plus youthful nostalgia.'' That was only one subject; folk can also be a village inside the city, a highway, a factory, a coast.

\subsection{2000s--2020s: Indie Folk and Global Localism}
\label{sec:folk-history-2000s}

In the 2000s, Anglo-American indie folk recombined banjo, accordion, vintage voices, and rock dynamics; Mumford \& Sons' 2009 ``Little Lion Man'' used banjo, stand-up bass, four-on-the-floor kick, and layered harmonies to make folk loud enough for a large stage.{\textsuperscript{[5]}} It was not ``returning to the past'' but using old instruments for new dynamics.

In China, indie folk spread through festivals, live houses, talent shows, and short video in the 2010s; Zhao Lei released ``Chengdu'' in 2016 and sang it on \emph{Singer} in 2017, placing second in the episode and turning Yulin Road, the little bar, and a private memory into a shared urban symbol.{\textsuperscript{[6]}} Localism in folk is not a few dialect lines; it is letting a city's details become the song's main character.

\subsection{Sources for the History}
\label{sec:folk-sources}

\begin{longtable}{p{0.06\textwidth}p{0.38\textwidth}p{0.46\textwidth}}
\toprule
No. & Fact & Source \\
\midrule
\endfirsthead
\toprule
No. & Fact & Source \\
\midrule
\endhead
1 & 1940s Woody Guthrie, Pete Seeger, and the American folk revival & Caixin, ``The Protester Who Never Left,'' \url{https://culture.caixin.com/2014-02-26/100643389.html} \\
2 & 1963 ``Blowin' in the Wind''; Dylan alone on guitar and harmonica & Bob Dylan official album page, \url{https://www.bobdylan.com/albums/freewheelin-bob-dylan/} \\
3 & 1971 \emph{Blue}; Joni Mitchell on dulcimer, James Taylor on acoustic guitar, Russ Kunkel on drums & Joni Mitchell official library, \url{https://jonimitchell.com/library/view.cfm?id=2545} \\
4 & 1994 \emph{Campus Folk I}; Lao Lang's ``My Deskmate'' & \emph{The Beijing News}, ``A History of Chinese Folk Music,'' \url{https://m.bjnews.com.cn/detail/157267174414834.html} \\
5 & 2009 ``Little Lion Man''; banjo, stand-up bass, four-on-the-floor kick, and vocal dynamics & Sound On Sound, ``Mumford \& Sons: Little Lion Man,'' \url{https://www.soundonsound.com/techniques/mumford-sons-little-lion-man} \\
6 & Zhao Lei's ``Chengdu'' released 2016; arrangement, production, and player credits & Douyin Baike, ``Chengdu (2016 song by Zhao Lei),'' \url{https://m.baike.com/wiki/%E6%88%90%E9%83%BD/19497957} \\
7 & 2017 \emph{Singer} performance of ``Chengdu,'' second place & CCTV, ``Zhao Lei and `Chengdu' Flood Social Media,'' \url{https://news.cctv.com/2017/02/08/ARTIsg0C6fU6tQYOM20KTuGo170208.shtml} \\
\bottomrule
\end{longtable}

\section{Arrangement: The Guitar Is a Second Voice, Not a Wall}
\label{sec:folk-craft}

\subsection{Write the Job Before the Instrument}
\label{sec:folk-instrumentation}

Folk can be two tracks: voice and acoustic guitar. It can be six to eight: voice, double, acoustic guitar, second guitar or banjo, bass, light drums, harmonica or flute, strings. More tracks do not make folk; each layer must ``speak.''

\begin{longtable}{p{0.15\textwidth}p{0.30\textwidth}p{0.27\textwidth}p{0.18\textwidth}}
\toprule
Layer & Common tracks/instruments & Main job & Common mistake \\
\midrule
\endfirsthead
\toprule
Layer & Common tracks/instruments & Main job & Common mistake \\
\midrule
\endhead
Voice & Lead, double, low-third harmony & Tell the story, breathe, carry accent and landing points & Doubling every word, so it sounds like a copied person \\
Guitar & Acoustic fingerpicking/strumming, second guitar & Set key, tempo, and leave gaps between phrases & Filling every beat, so the voice has no room to breathe \\
Low end & Bass, stand-up bass, low kick & Hold the root and give the chorus weight & Full drum kit in the verse, so the chorus has no contrast \\
Colour & Harmonica, banjo, accordion, flute, strings & Mark section borders and place & Colour from start to finish, becoming background noise \\
Space & Room tone, close noise, reverb & Make the listener feel ``in a place'' & Faking sincerity with a lo-fi filter \\
\bottomrule
\end{longtable}

A reproducible start is six tracks: Vocal (audio), Guitar (audio), Guitar Double (audio, chorus only), Bass (audio or MIDI), Percussion (audio/MIDI, can mute in verses), Pad/Harmony (MIDI, added last). Make voice and one guitar stand first; if those two cannot stand, more tracks only hide the problem.

\subsection{Texture: Fingerpicking, Strumming, Travis Picking}
\label{sec:folk-texture}

Three common acoustic-guitar textures:

\begin{itemize}
\item \textbf{Strumming}: good for choruses and for ``pushing the words out''; the danger is strumming the verse too, which is loud and leaves no landing place.
\item \textbf{Arpeggios}: good for narrative verses; the danger is repeating the same bar mechanically, so change the pattern or leave half a beat at phrase ends.
\item \textbf{Travis picking / alternating bass}: thumb holds the bass, index and middle fingers pick the melody; intimate, like one person in a room, and easiest to overfill.
\end{itemize}

Ask first: does this line need the listener to hear a breath? If yes, let the guitar stop for half a beat at the phrase end; if no, let it keep walking and push the voice forward. Folk's ``breath'' is not low volume; it is \textbf{space between phrases}.

\subsection{Harmony: A Simple Cycle With Clear Landings}
\label{sec:folk-harmony}

I--IV--V--vi is folk's most common grammar, but simple does not mean careless. Three practical rules:

\begin{enumerate}
\item \textbf{Give the bass a direction}: do not return to the tonic every bar; use inversions, suspensions, and open chords to move the bass two steps.
\item \textbf{Change one thing in the chorus}: swap vi for I, change a major third to a minor third, or stretch the last beat; do not change chords, tempo, and texture at once.
\item \textbf{Open tunings are for timbre, not showing off}: DADGAD or Open D gives longer resonance and a lower root; check the singer's range first, and do not tune the song into an uncomfortable key.
\end{enumerate}

\subsection{Dynamics: The Chorus Is ``One More Person,'' Not a Fader}
\label{sec:folk-dynamic}

Folk's most effective chorus lift is usually not a volume push but ``one more person/one more instrument'':

\begin{longtable}{p{0.18\textwidth}p{0.36\textwidth}p{0.38\textwidth}}
\toprule
Section & What can enter & What not to do \\
\midrule
\endfirsthead
\toprule
Section & What can enter & What not to do \\
\midrule
\endhead
Verse & Voice plus acoustic guitar; brushes or tambourine if needed & No full drum kit; do not strum every beat \\
Pre-chorus & Bass enters; second guitar enters; harmony only at phrase ends & Do not lay harmony from start to finish \\
Chorus & Double/low third; light kick; one line of harmonica or flute & Do not push every instrument to maximum at once \\
Bridge & Remove drums and bass; voice plus guitar only & Do not add an electronic drop for ``change'' \\
Outro & Return to the verse motive; leave space & Do not hide a missing ending under long reverb \\
\bottomrule
\end{longtable}

\subsection{The Yeban Skeleton: Start With Six Tracks}
\label{sec:folk-yeban}

\begin{longtable}{p{0.13\textwidth}p{0.20\textwidth}p{0.55\textwidth}}
\toprule
Track & Type & What to write first \\
\midrule
\endfirsthead
\toprule
Track & Type & What to write first \\
\midrule
\endhead
Vocal & Audio & An eight-bar verse with half a beat to a beat of space at each phrase end; do not tune first, check that it sounds like talking \\
Guitar & Audio & Arpeggios or Travis picking; leave a gap at each phrase end, do not fill every beat \\
Guitar Double & Audio & Chorus only; 3--6 dB below the main guitar, not a second lead \\
Bass & Audio/MIDI & Root only in the verse; passing tones in the chorus; do not fight the kick for the same beat \\
Percussion & Audio/MIDI & Can mute in the verse; brushes or light kick in the chorus; no full kit \\
Pad/Harmony & MIDI & Add last; long notes or phrase-end harmony only; do not cover the guitar's midrange \\
\bottomrule
\end{longtable}

Order: voice → guitar → bass → chorus double → light drums → colour. After each layer, solo it and ask: ``Which line does this help?'' If you cannot answer, remove it.

\section{Four Songs, Four Ways of Doing Less}
\label{sec:folk-cases}

\subsection{\emph{Blowin' in the Wind}: One Person Is a Band}
\label{sec:folk-case-dylan}

In 1963, Bob Dylan sang ``Blowin' in the Wind'' on \emph{The Freewheelin' Bob Dylan}. The official page is clear: on this track Dylan accompanies himself alone on acoustic guitar and harmonica---no drums, no band.{\textsuperscript{[2]}} The song moves by repeated questions; each verse pushes the answer further away and leaves the judgment to the listener.

Its arrangement intelligence is ``not resolving'': the guitar keeps a steady cycle, the harmonica acts like a second speaker at phrase ends, and the voice does not aim for prettiness but for ``this must be said.'' Nothing in the arrangement shows off, yet a whole era sang along.

\textbf{One thing to steal}: if you have only one guitar, do not rush to add backing. Write a cycle that can repeat, then decide where the harmonica or second voice appears at a phrase end. In a one-person song, space is the band.

\subsection{\emph{A Case of You}: Timbre Speaks Louder Than Track Count}
\label{sec:folk-case-joni}

Joni Mitchell's 1971 ``A Case of You'' is on \emph{Blue}. The official library notes a restrained accompaniment: Joni plays Appalachian dulcimer, James Taylor plays acoustic guitar, Russ Kunkel plays drums, and the voice is placed in front.{\textsuperscript{[3]}} The dulcimer's metallic strings and the guitar's warm grain sit side by side, like two different ``I''s in conversation.

It does not turn love into a big chorus; timbre carries the emotion. The dulcimer's brightness, the guitar's softness, the drums' lightness---three instruments each say one sentence. The chorus does not add tracks; it lets the voice slide from murmur to high register.

\textbf{One thing to steal}: find one timbre outside the standard band---dulcimer, harmonica, accordion, even a guitar tuned very low. Let it appear two or three times, not from start to finish; that is more memorable.

\subsection{\emph{Little Lion Man}: Build the Chorus by Letting People In}
\label{sec:folk-case-mumford}

Mumford \& Sons' 2009 ``Little Lion Man'' was their debut single, on \emph{Sigh No More}. The Sound On Sound production interview notes that the album was mostly recorded live in the studio, with drums overdubbed later; the drums are mostly four-on-the-floor kick and floor tom, with occasional snare; the instrumentation includes stand-up bass, banjo, piano, Marcus Mumford's voice and acoustic guitar, and lots of backing vocals.{\textsuperscript{[5]}}

Its chorus does not ``turn up the volume''; ``more people enter.'' Acoustic guitar first, then banjo, piano, bass, kick, and finally harmonies lift the room. The banjo is assertive, but the production deliberately makes it ``cut through without becoming a banjo record''---that is dynamic allocation.

\textbf{One thing to steal}: do not put every instrument in before the chorus. Let each enter on a different beat of the chorus; the listener feels the song growing, not a fader being pushed.

\subsection{Zhao Lei's ``Chengdu'': Make a City the Second Lead}
\label{sec:folk-case-chengdu}

Zhao Lei released ``Chengdu'' in 2016; it was included on the album \emph{Unable to Grow Up} in December 2016. The credits list Zhao Lei and Xizi for arrangement/production, Zhang Ling on bass, Beibei on drums, Liu Sen on piano, Zhu Zi on flute, Zhao Lei and Xizi on acoustic guitar, plus string arrangement.{\textsuperscript{[6]}} In 2017 he sang it on \emph{Singer}, placed second in the episode, and pushed indie folk into the mainstream.{\textsuperscript{[7]}}

The arrangement does not turn Chengdu into a tourism video: the verse is a close-up of acoustic guitar and voice, bass and drums are light, and flute and strings appear only at section borders, like a sound from around the corner. It lets concrete names---Yulin Road, the little bar---carry the narrative instead of stacking adjectives.

\textbf{One thing to steal}: do not write ``this city is beautiful''; write three objects/places/times only locals know, and give just one of them one timbre (flute, accordion, a piece of room sound). Do not give all three everything.

\section{Three Exercises: Take Folk Apart at Yeban}
\label{sec:folk-practice}

\begin{exercise}
\textbf{One lyric, three guitar textures.} Take an eight-bar verse and make A arpeggios, B strumming, C Travis picking. Record the vocal on all three and listen blind: which sounds most like ``talking,'' which most like ``performing,'' which makes the chorus easiest to enter. A is intimate but can go flat; B pushes but can get loud; C has texture but easily overfills. There is no right answer, only trade-offs.
\end{exercise}

\begin{exercise}
\textbf{The removal test.} Make a ``full band'' demo (voice, guitar, bass, drums, pad, harmony), then strip it to voice plus guitar, then add back one instrument at a time. Each time ask: which line did it help? If it only makes the song ``fuller,'' do not add it. The final three to five tracks usually sound more like folk than eight.
\end{exercise}

\begin{exercise}
\textbf{One city, three concrete objects.} Write eight lines without big words such as ``beautiful,'' ``lonely,'' or ``far away''; use only three specific places/objects/times (a shop, a bus, a rain). Set it to a I--IV--V--vi cycle; verse with guitar only, chorus with one colour instrument (flute/harmonica/accordion). Make sure the song stands a cappella before discussing arrangement.
\end{exercise}

\section{Further Reading}
\label{sec:folk-reading}

\begin{sloppypar}
\begin{itemize}
\item \textbf{Bob Dylan official page for \emph{The Freewheelin' Bob Dylan}}: \url{https://www.bobdylan.com/albums/freewheelin-bob-dylan/}. Check the recording, guitar/harmonica scoring, and album context of ``Blowin' in the Wind.''
\item \textbf{Joni Mitchell official page for \emph{Blue}}: \url{https://jonimitchell.com/music/album.cfm?id=5}. See where ``A Case of You'' sits on the album and how dulcimer, acoustic guitar, and drums are kept restrained.
\item \textbf{Sound On Sound, ``Mumford \& Sons: Little Lion Man''}: \url{https://www.soundonsound.com/techniques/mumford-sons-little-lion-man}. Use the production interview to understand banjo, stand-up bass, four-on-the-floor kick, and vocal dynamics.
\item \textbf{CCTV, ``Zhao Lei and `Chengdu' Flood Social Media''}: \url{https://news.cctv.com/2017/02/08/ARTIsg0C6fU6tQYOM20KTuGo170208.shtml}. Check the song's spread after \emph{Singer} and the mainstreaming of indie folk.
\end{itemize}
\end{sloppypar}

These links were checked on November 27, 2026 and were reachable. If a regional restriction blocks access, search the title in the platform's official account.

\section{Chapter Takeaways}
\label{sec:folk-takeaways}

\begin{itemize}
\item Folk is not ``pop with a smaller budget''; it puts a specific person, place, and moment back into the song.
\item The voice is the lead, the guitar is a second voice; every guitar beat must leave the voice room to breathe.
\item Lift the chorus by adding a person or an instrument, not by pushing the fader.
\item Timbre speaks louder than track count: one instrument outside the standard band, appearing two or three times, is enough to remember.
\item City narrative relies on concrete objects, not big words; give one object one timbre.
\item Make the a cappella and one guitar stand first; if they cannot, more tracks only hide the problem.
\end{itemize}

\begin{helpmapping}
\begin{longtable}{p{0.45\textwidth}p{0.45\textwidth}}
\toprule
Section & Help anchor \\
\midrule
What folk is & \texttt{arrangement-genre-folk-what} \\
The route & \texttt{arrangement-genre-folk-history} \\
Arrangement & \texttt{arrangement-genre-folk-craft} \\
Four songs & \texttt{arrangement-genre-folk-cases} \\
Yeban skeleton & \texttt{arrangement-genre-folk-yeban} \\
Exercises & \texttt{arrangement-genre-folk-practice} \\
Further reading & \texttt{arrangement-genre-folk-reading} \\
\bottomrule
\end{longtable}
\end{helpmapping}
